\clearpage
\SourceLeaf{145}
\begin{center}
\textsc{Gauss} an \textsc{Schumacher}. Göttingen, 3. December 1831.
\end{center}

\begin{center}
. . . . . . . . . . .
\end{center}

Die von Ihnen erwähnte Stelle in \textsc{Zachs} A.\,G.\,E. ist mir wohl bekannt; die Anwendung der M. der kl. Q., deren dort Erwähnung geschieht, betrifft einen früher in derselben Zeitschrift abgedruckten Auszug aus \textsc{Ulugh Beighs} Zeitgleichungs-Tafel, die zu manchen ganz curiosen Resultaten geführt hatte. Diese Resultate hatte ich \textsc{Zach} mitgetheilt mit der Bemerkung, dass ich dabei eine mir eigenthümliche seit Jahren gebrauchte Methode benutzt habe, Grössen, die zufällige Fehler involviren, auf eine willkürfreie consequente Art zu combiniren, ohne ihm jedoch das Wesen der Methode selbst mitzutheilen. Ich glaube Ihnen schon einmal geschrieben zu haben, dass ich auf keinen Fall diese Stelle, worin die Methode zum erstenmale öffentlich angedeutet ist, releviren werde, auch nicht wünsche, dass einer meiner Freunde mit meiner Zustimmung es thue. Diess hiesse \textls{anerkennen, als bedürfe} meine Anzeige (Theoria Motus Corporum Coelestium), dass ich seit 1794 diese Methode vielfach gebraucht habe, einer Rechtfertigung, und dazu werde ich mich nie verstehen. Als \textsc{Olbers} attestirte\HistoricalNote{[*)]}{[*)}{In seiner Abhandlung: Über den veränderlichen Stern im Halse des Schwans; v.\,\textsc{Lindenau} und \textsc{Bohnenberger}, Zeitschrift für Astronomie. Band II. Seite 192. September-October 1816.]}, dass ich ihm 1802\HistoricalNote{[**)]}{[**)}{Sollte heissen »1803«.]} die ganze Methode mitgetheilt habe, war diess zwar gut gemeint; hätte er mich aber vorher gefragt, so würde ich es hautement gemissbilligt haben.

\begin{center}
. . . . . . . . . . . . . . .
\end{center}

\begin{center}
\textsc{Gauss} an \textsc{Olbers}. Braunschweig, 30. Juli 1806.
\end{center}

\begin{center}
. . . . . . . . . . .
\end{center}

Hr.\ \textsc{von Zach} schreibt mir noch, dass Sie Sich zur Recension von \textsc{Legendres} Werk über die Kometenbahnen erboten hätten. Mit Vergnügen werde ich Ihnen also das mir von Hrn.\ \textsc{von Zach} zugeschickte Exemplar nach Bremen
